\section{Setup: agents, markets and instruments}\label{sec:mkt:setup}

Part~\ref{part:sp} solved the planner's problem under two closures for the material intensity of new
capital goods. This part states the market economy that corresponds to it, and asks what a government
must do to make a competitive equilibrium reproduce the planner's allocation. Tastes and technologies
are unchanged: everything in Section~\ref{sec:sp:setup:primitives}, the process-level accounting of
Section~\ref{sec:sp:setup:materialaccounting}, the embodied-material
stock~\eqref{eq:sp:setup:Phi-motion} and the ledger~\eqref{eq:sp:setup:ledger} carry over verbatim.
What is new here is institutional: who owns what, what is traded at what price, and which of the
model's shadow prices has no private counterpart.


\subsection{What is decentralised and what is not}\label{sec:mkt:setup:what}

One point of method is worth stating before anything else, because it governs the whole part. The
materials accounting closes \emph{identically} in the two economies. The
ledger~\eqref{eq:sp:setup:ledger} determines the waste flow $W$, and~\eqref{eq:sp:setup:Omega}
or~\eqref{eq:sp:setup:Omega-phi} determines the intensity $\Omega$, in the market exactly as in the
planner's problem, whatever policy is in force and whether or not the allocation is optimal.
Conservation of mass is physics, not an institution, and there is no sense in which a planner and a
market face different versions of it. What differs between the two economies is only what private
agents \emph{internalise}, and that difference is the entire subject of this part.

A second structural fact organises the presentation, and it is worth stating equally early because it
is what makes this part shorter than Part~\ref{part:sp}.

\begin{proposition}[The market structure does not depend on the closure]\label{prop:mkt:common}
The agents, markets, prices and optimality conditions of the decentralised economy are the same under
Version~B and Version~A. The two closures differ only in (i) how the accounting
determines $W$ and $\Omega$ given the allocation, which no agent chooses, and (ii) the values at which
the policy instruments must be set for the equilibrium to be optimal.
\end{proposition}

The reason is that neither $\Omega$ nor any relative intensity $\omega^j$ appears in any private
agent's problem below. Firms see physical technologies and market prices; the accounting layer is
invisible to them. It enters their decisions only through the instruments the government chooses, and
that is where the closure matters. Accordingly the agents' problems and the equilibrium concept are
stated once, here; Sections~\ref{sec:mkt:A} and~\ref{sec:mkt:B} then set the instruments under each
closure and verify implementation.


\subsection{Agents and markets}\label{sec:mkt:setup:agents}

The economy contains five private decision units and a government. The final good is the numéraire.

\smalltitle{The final-good firm.} A competitive firm operates the technology~\eqref{eq:sp:setup:output},
renting capital $K^Y$ and buying raw material $R$. It takes the pollution stock $P$ parametrically.

\smalltitle{The mine.} A competitive firm owns the reserve $S$ and the record of cumulative discovery
$X$, operates the cost functions~\eqref{eq:sp:setup:extraction-cost}--\eqref{eq:sp:setup:discovery-cost},
extracts $N$ and explores $D$, and sells virgin material. It is a dynamic agent: it holds two stocks
and therefore carries two private shadow prices.

\smalltitle{The recycling sector.} Competitive firms buy treated waste $T$ from the waste-management
sector, rent capital $K^R$, and operate the constant-returns
technology~\eqref{eq:sp:setup:recycling-crs}, selling secondary material into the same market as the
mine. Virgin and secondary material are physically identical, so there is one material price.

\smalltitle{The waste-management sector.} Competitive firms operate the collection and treatment
technology~\eqref{eq:sp:setup:waste-cost} under a public service obligation: \emph{all} waste generated
must be collected. They choose the treated share $\varpi$, sell treated waste to recyclers, and release
the remainder to the environment. That the obligation is enforced --- that there is no unrecorded
disposal --- is an institutional assumption, and a substantive one; see~\ref{op:mkt:dumping}.

\smalltitle{The household.} A representative dynasty with preferences~\eqref{eq:sp:setup:utility} owns
the capital stock, rents it to the two capital-using sectors, receives the profits of the firms it
owns, pays taxes, receives a transfer, and consumes. It takes $P$ parametrically.

\smalltitle{The government.} Sets the instruments of Section~\ref{sec:mkt:setup:instruments} and
balances its budget period by period with a lump-sum transfer $T$, which is unrestricted in sign.

Four prices are determined in markets, and are listed with the two private shadow prices in
Table~\ref{tab:mkt:prices}. All are measured in units of the final good.

\begin{table}[H]
\centering
\caption{Prices in the decentralised economy.}
\label{tab:mkt:prices}
\setlength{\tabcolsep}{5pt}
\renewcommand{\arraystretch}{1.05}
\begin{tabular}{@{}p{0.10\textwidth}p{0.42\textwidth}p{0.40\textwidth}@{}}
\toprule
Price & Market & Paid by / to \\
\midrule
$P^R$ & raw material, virgin and secondary  & final-good firm to mine and recyclers \\
$P^T$ & treated waste                       & recyclers to waste-management sector \\
$P^K$ & rental of capital services          & final-good firm and recyclers to household \\
$Q$   & purchase price of new capital goods & household, per unit of gross formation \\
\midrule
$P^S$ & shadow price of the reserve         & internal to the mine \\
$P^X$ & shadow price of cumulative discovery & internal to the mine \\
\bottomrule
\end{tabular}
\end{table}

There is no market for waste \emph{generation}: nobody buys the right to create a tonne of waste, and
the collection sector does not levy a gate fee on those who hand waste over. Whether it should is not a
detail of presentation --- it is the substantive policy question of this part, taken up in
Sections~\ref{sec:mkt:A:gate} and~\ref{sec:mkt:B:nogate}.


\subsection{Policy instruments}\label{sec:mkt:setup:instruments}

Table~\ref{tab:mkt:instruments} lists the instruments. All may be functions of time, all are
unrestricted in sign, and all are taken as given by private agents. Two are not free: the recycling
credit $s^R$ is tied to the emission tax by construction, and the collection payment $s^W$ is whatever
the competitive tender requires for the waste sector to break even.

\begin{table}[H]
\centering
\caption{Policy instruments. The last column anticipates the Pigouvian values derived in
Sections~\ref{sec:mkt:A} and~\ref{sec:mkt:B}, stated for Version~A: $\zeta$ is the intensity
multiplier of~\eqref{eq:sp:A:zeta} and $\Delta=\zeta\Omega$ the composition wedge. Under Version~B the
deposit takes the same form with the mass-constraint multiplier $\eta$ of~\eqref{eq:sp:B:foc:cs} in
place of $\zeta$, and $\Delta\equiv0$, so the entire lower block vanishes.}
\label{tab:mkt:instruments}
\setlength{\tabcolsep}{4pt}
\renewcommand{\arraystretch}{1.05}
\begin{tabular}{@{}p{0.09\textwidth}p{0.20\textwidth}p{0.38\textwidth}p{0.25\textwidth}@{}}
\toprule
& Base & Levied on & Optimal value \\
\midrule
$\tau^R$      & $R$                  & material entering production, either source & $p^W-\zeta$ \\
$s^W$         & $W$                  & waste collected (residual, from zero profit) & $p^W$ \\
$\tau^{N,S}$  & $N$                  & tonnage extracted (displaced overburden) & $\Omega^{N,S}p^W$ \\
$\tau^{D,S}$  & $D$                  & tonnage discovered (displaced overburden) & $\Omega^{D,S}p^W$ \\
$\tau^E$      & $\left(1-\varpi+d^W\varpi\right)W$ & waste released to the environment & $p^P$ \\
$s^R$         & $R^R$                & secondary material recovered (credit against $\tau^E$) & $d^W\tau^E$ \\
$\tau^K$      & $G=I+\delta K$       & \emph{gross} capital formation & see~\eqref{eq:mkt:A:tauK}, \eqref{eq:mkt:B:tauK} \\
\midrule
$\tau^Y$      & $Y$                  & output of the final good & $\Delta\omega^Y$ \\
$\tau^C$      & $C$                  & consumption & $\Delta\omega^C$ \\
$\tau^N$      & $C^N$                & extraction expenditure & $\Delta\omega^N$ \\
$\tau^D$      & $C^D$                & exploration expenditure & $\Delta\omega^D$ \\
$\tau^W$      & $C^W$                & collection and treatment expenditure & $\Delta\omega^W$ \\
\bottomrule
\end{tabular}
\end{table}

The division of the table is the division of the part. The upper block prices \emph{mass and
emissions}: it is present under both closures and its members are the instruments a materials-balance
model must have. The lower block prices the \emph{composition of final-good spending}: every member of
it is proportional to $\Delta$, and by Proposition~\ref{prop:sp:B:irrelevance} $\Delta\equiv0$ under
Version~B, so the entire block vanishes there. An economy of Version~B type needs no tax on output and
no tax on consumption; an economy of Version~A type needs one on every use of the final good. This is
the market counterpart of the comparison in Section~\ref{sec:sp:B:compare} and, unlike most such
comparisons, it is directly observable in the policy menu.

Notation deserves one warning. The five instruments in the lower block are indexed by the activity
whose \emph{expenditure} they fall on, so $\tau^N$ is an ad valorem tax on extraction \emph{effort}
$C^N(N,S)$, whereas $\tau^{N,S}$ is a specific tax on the \emph{tonnage} $N$; the superscript $S$
marks the nature-attributable base, as in $\Omega^{N,S}$. The two are different instruments with
different bases and, as it turns out, different fates across the two closures.


\subsection{The agents' problems}\label{sec:mkt:setup:problems}

Throughout, $r$ denotes the market rate of interest, defined in~\eqref{eq:mkt:setup:household-euler}
below, and firms with dynamic problems discount at it.

\smalltitle{The final-good firm.} The firm solves
\begin{align}
\max_{K^Y,R}\;\left(1-\tau^Y\right)F\left(K^Y,R,P\right)-P^KK^Y-\left(P^R+\tau^R\right)R,
\label{eq:mkt:setup:goodsfirm}
\end{align}
giving
\begin{align}
\left(1-\tau^Y\right)F_K &= P^K,
&
\left(1-\tau^Y\right)F_R &= P^R+\tau^R .
\label{eq:mkt:setup:foc:goodsfirm}
\end{align}
The firm buys mass and pays $\tau^R$ on it; it does not pay for the waste that mass will become,
because it never handles that waste. Whether the two are the same charge in disguise is the question
Sections~\ref{sec:mkt:A:reading} and~\ref{sec:mkt:B:reading} answer.

\smalltitle{The recycling sector.} A recycler buys treated waste and rents capital,
\begin{align}
\max_{T,K^R}\;\left(P^R+s^R\right)\mathcal R\left(T,K^R\right)-P^TT-P^KK^R,
\label{eq:mkt:setup:recycler}
\end{align}
and, using the derivatives~\eqref{eq:sp:setup:recycling-derivatives} of the constant-returns
technology,
\begin{align}
\alpha\left(P^R+s^R\right) &= P^T,
&
a'\!\left(x\right)\left(P^R+s^R\right) &= P^K,
&
x&=\frac{K^R}{T}.
\label{eq:mkt:setup:foc:recycler}
\end{align}
Because $\mathcal R$ is homogeneous of degree one, $\mathcal R=\alpha T+a'K^R$
by~\eqref{eq:sp:setup:recycling-derivatives}, and the two conditions
in~\eqref{eq:mkt:setup:foc:recycler} imply that the sector's revenue is exhausted by its factor
payments: recycling earns exactly zero profit at every date, on and off the optimal path. The object
\begin{align}
B &\equiv P^R+s^R
\label{eq:mkt:setup:Bvalue}
\end{align}
--- the price a recycler receives per tonne of secondary material, inclusive of the credit --- is the
sector's sufficient statistic, and~\eqref{eq:mkt:setup:foc:recycler} says that treated waste is worth
the marginal recycling yield times that price. The notation is deliberate: $B$ is the market
counterpart of the planner's~\eqref{eq:sp:B:Bvalue} and~\eqref{eq:sp:A:Bvalue}, and the implementation
propositions below show that the two coincide.

\smalltitle{The mine.} The mine maximises the present value of its dividends,
\begin{align}
\max_{N,D}\;\int_0^\infty\Big[P^RN-\left(1+\tau^N\right)C^N\!\left(N,S\right)
-\left(1+\tau^D\right)C^D\!\left(D,X\right)-\tau^{N,S}N-\tau^{D,S}D\Big]
e^{-\int_0^tr\,ds}dt,
\label{eq:mkt:setup:mine}
\end{align}
subject to $\dot S=D-N$ and $\dot X=D$, with $S_0,X_0$ given. Writing $P^S$ and $P^X$ for the
current-value costates on the two stocks,
\begin{align}
P^R &= \left(1+\tau^N\right)C^N_N+P^S+\tau^{N,S},
&
P^S+P^X &= \left(1+\tau^D\right)C^D_D+\tau^{D,S},
\label{eq:mkt:setup:foc:mine}
\\
\dot P^S &= r\,P^S+\left(1+\tau^N\right)C^N_S,
&
\dot P^X &= r\,P^X+\left(1+\tau^D\right)C^D_X .
\label{eq:mkt:setup:costate:mine}
\end{align}
Both stock effects are internalised, because the mine owns both stocks: the depletion effect $C^N_S$
because it will pay the higher extraction cost itself, and the discovery-depletion effect $C^D_X$
because it will pay the higher exploration cost itself. There is accordingly no instrument in
Table~\ref{tab:mkt:instruments} attached to either. This is worth marking, since the resource-economics
literature often supplies one: here the reserve and the exploration margin are private, and the only
reason the mine's decisions are distorted at all is the mass its activity moves.

\smalltitle{The waste-management sector.} A collector takes the waste flow $W$ as given --- it is
obliged to collect it --- and chooses the treated share:
\begin{align}
\Pi^W &= \Big[P^T\varpi+s^W-\left(1+\tau^W\right)c^W\!\left(\varpi\right)
-\tau^E\left(1-\varpi+d^W\varpi\right)\Big]W,
&
c^W\!\left(\varpi\right)&=c^c+c^T\!\left(\varpi\right),
\label{eq:mkt:setup:wastefirm}
\end{align}
with first-order condition and zero-profit condition
\begin{align}
P^T+\left(1-d^W\right)\tau^E &= \left(1+\tau^W\right)\frac{dc^T}{d\varpi},
\label{eq:mkt:setup:foc:varpi}
\\
s^W &= \left(1+\tau^W\right)c^W\!\left(\varpi\right)
+\tau^E\left(1-\varpi+d^W\varpi\right)-P^T\varpi .
\label{eq:mkt:setup:zeroprofit}
\end{align}
The tender drives profit to zero, so~\eqref{eq:mkt:setup:zeroprofit} determines the collection payment
rather than the government choosing it. If treated waste is valuable enough, $s^W$ turns negative and
the ``payment'' becomes a licence fee for the right to collect. Note that no waste-intensity parameter
appears in~\eqref{eq:mkt:setup:wastefirm}: the collector sees the physical technology and nothing else.

\smalltitle{The household.} The household holds the capital stock, buys gross formation $G=I+\delta K$
at the price $Q\equiv1+\tau^K$, and consumes at the tax-inclusive price $1+\tau^C$:
\begin{align}
\max_{C}\;\int_0^\infty\left[u(C)-v(P)\right]e^{-\rho t}dt
\qquad\text{s.t.}\qquad
Q\left(\dot K+\delta K\right) &= P^KK+\Pi+T-\left(1+\tau^C\right)C,
\label{eq:mkt:setup:household}
\end{align}
where $\Pi$ collects the dividends of the firms it owns, $P$ is taken as given, $K_0$ is given, and a
no-Ponzi condition holds. Writing $\lambda$ for the household's marginal utility of a unit of the final
good and $\mu=\lambda Q$ for the costate on $K$,
\begin{align}
u'\!\left(C\right) &= \lambda\left(1+\tau^C\right),
&
r &\equiv \rho-\frac{\dot\lambda}{\lambda} \;=\; \frac{P^K+\dot Q}{Q}-\delta .
\label{eq:mkt:setup:household-euler}
\end{align}
The second relation defines the market interest rate and simultaneously states the household's
portfolio condition: the return on holding capital, which is the rental net of depreciation charged at
the price of capital, plus the capital gain, must equal the rate at which the household discounts.
Since $\tau^K$ is levied on gross formation, replacement of depreciated capital is priced at $Q$
exactly as net accumulation is. That is what the accounting requires --- both flows draw material on
identical terms, by~\eqref{eq:sp:setup:Phi-motion} --- and it is the reason the base is gross rather
than net.

\smalltitle{The government.} The transfer balances the budget at every date,
\begin{align}
T &= \tau^YY+\tau^CC+\tau^NC^N+\tau^DC^D+\tau^WC^W
+\tau^RR+\tau^{N,S}N+\tau^{D,S}D+\tau^KG
\notag\\
&\phantom{{}={}}
+\tau^E\left(1-\varpi+d^W\varpi\right)W-s^RR^R-s^WW .
\label{eq:mkt:setup:govbudget}
\end{align}


\subsection{Equilibrium}\label{sec:mkt:setup:equilibrium}

\begin{definition}[Competitive equilibrium]\label{def:mkt:equilibrium}
Given instrument paths and initial stocks $K_0,S_0,X_0,P_0,M^K_0$, a competitive equilibrium is a path
of quantities $\left\{C,I,N,D,\varpi,K^Y,K^R,T,R,R^R,W,Y\right\}$, prices
$\left\{P^R,P^T,P^K,Q,r\right\}$, private shadow prices $\left\{P^S,P^X\right\}$ and transfers $T$ such
that each agent solves its problem in Section~\ref{sec:mkt:setup:problems}; the markets for material,
treated waste and capital services clear,
\begin{align}
R &= N+R^R,
&
T &= \varpi W,
&
K &= K^Y+K^R;
\label{eq:mkt:setup:clearing}
\end{align}
the government budget~\eqref{eq:mkt:setup:govbudget} holds; the stocks evolve
by~\eqref{eq:sp:setup:reserves-motion}, \eqref{eq:sp:setup:pollution-motion}
and~\eqref{eq:sp:setup:Phi-motion} with the emission flow~\eqref{eq:sp:setup:Xi}; and the materials
accounting holds, that is $W$ satisfies the ledger~\eqref{eq:sp:setup:ledger} and $\Omega$ is given by
the closure in force.
\end{definition}

Two features of the definition are worth marking. First, $W$ is not chosen by anybody. It is
determined by the ledger given the allocation, exactly as in Part~\ref{part:sp}, where it was carried
as a control only to keep its multiplier visible. Second, the last clause is the only place the two
closures enter, and it does not touch any agent's problem --- which is
Proposition~\ref{prop:mkt:common}.

\smalltitle{Walras' law.} The instrument set and the transfer rule are mutually consistent, in the
following sense.

\begin{proposition}[The household's budget is the economy's resource constraint]
\label{prop:mkt:walras}
In any competitive equilibrium, the household's budget
constraint~\eqref{eq:mkt:setup:household} together with firm profits, the zero-profit conditions of the
recycling and waste-management sectors, market clearing~\eqref{eq:mkt:setup:clearing} and the
government budget~\eqref{eq:mkt:setup:govbudget} implies the goods
constraint~\eqref{eq:sp:setup:resource-constraint},
\begin{align}
Y &= C+I+\delta K+C^N+C^D+C^W .
\label{eq:mkt:setup:resource}
\end{align}
This holds for arbitrary instrument values, not only at the optimum.
\end{proposition}

\begin{proof}
Substitute the dividends of the final-good firm and the mine, and the zero-profit conditions of the
recycling and waste-management sectors, into~\eqref{eq:mkt:setup:household}. Capital rentals cancel by
$K=K^Y+K^R$. Material payments cancel by $R=N+R^R$, leaving $-\tau^RR+s^RR^R$. Payments for treated
waste cancel by $T=\varpi W$, leaving, through~\eqref{eq:mkt:setup:zeroprofit},
$-\left(1+\tau^W\right)C^W-\tau^E\left(1-\varpi+d^W\varpi\right)W+s^WW$. Every remaining tax and
subsidy term then appears once with each sign, once in the household's budget and once
in~\eqref{eq:mkt:setup:govbudget}, and cancels; for $\tau^K$ this is the statement that the household
pays $\left(1+\tau^K\right)G$ and is returned $\tau^KG$. What is left is
$Y-C-I-\delta K-C^N-C^D-C^W=0$.
\end{proof}

The proposition is the check that the instrument bases are the right ones. It is not automatic: a
$\tau^K$ levied on net investment, or an emission tax whose base did not coincide with the physical
release, would fail it.


\subsection{The externalities, and the unregulated equilibrium}\label{sec:mkt:setup:laissezfaire}

Private agents fail to internalise exactly two things, and it is worth naming them precisely because
the instrument menu is long enough to obscure how few underlying distortions there are.

\begin{enumerate}[label=\textbf{X\arabic*.},ref=X\arabic*,leftmargin=2.6em,itemsep=0.3em]
\item\label{ext:mkt:pollution} \textbf{The pollution stock.} Household and firms take $P$ as given,
so neither the disutility $v'(P)$ nor the output loss $F_P$ is priced. There is no private counterpart
to $p^P$.
\item\label{ext:mkt:ledger} \textbf{The materials ledger.} No agent faces~\eqref{eq:sp:setup:ledger}.
The final-good firm draws mass without paying for the waste it becomes; the household stores mass in
durables without being credited for the deferral; the mine displaces overburden without paying for it.
There is no private counterpart to $p^W$ or to the embodied liability $p^M$.
\end{enumerate}

Everything else is internal. The reserve rent and the exploration externality are the mine's own
problem; the recycling technology is priced by $P^T$ and $P^K$; the collection technology is priced by
the tender. Under Version~A a third margin --- the effect of the composition of spending on how much
mass is stored rather than shed --- is likewise unpriced, but it is a facet of~\ref{ext:mkt:ledger}
rather than a separate failure, and it is why the lower block of Table~\ref{tab:mkt:instruments}
exists.

\smalltitle{The unregulated equilibrium.} Setting every instrument to zero
in~\eqref{eq:mkt:setup:foc:goodsfirm}--\eqref{eq:mkt:setup:household-euler} gives a benchmark worth
recording, since it is the natural status quo for a counterfactual exercise. It is not a no-recycling
economy. With $\tau^E=0$ and $s^W=0$ the collection sector must cover its costs from the sale of
treated waste alone, so~\eqref{eq:mkt:setup:foc:varpi} and~\eqref{eq:mkt:setup:zeroprofit} become
\begin{align}
\frac{dc^T}{d\varpi} &= P^T=\alpha P^R,
&
\varpi\frac{dc^T}{d\varpi}-c^T\!\left(\varpi\right) &= c^c,
\label{eq:mkt:setup:lf-waste}
\end{align}
the second obtained by eliminating $P^T$ between the first-order condition and zero profit. The
unregulated economy is therefore not a no-recycling economy --- secondary material is worth something,
and recyclers bid for treated waste --- but the way it recycles is peculiar.

\begin{proposition}[The unregulated treatment share is a technological constant]
\label{prop:mkt:lf-varpi}
With every instrument set to zero, the treated share $\varpi$ and the price of treated waste $P^T$ are
determined by the collection technology alone, from~\eqref{eq:mkt:setup:lf-waste}, independently of
resource scarcity, damages and every price in the economy. The left-hand side of the second condition
is zero at $\varpi=0$ and strictly increasing, with derivative $\varpi\,d^2c^T/d\varpi^2>0$, so the
solution is unique; it is increasing in the pure collection cost $c^c$.
\end{proposition}

The economics is that the collector's only revenue is the sale of treated waste, so it treats up to the
point where that revenue covers the whole of collection and treatment --- a condition on costs, not on
values. The consequence matters for the mechanism of the paper. Under the optimal policy, scarcity
raises circularity through \emph{both} margins of the recycling block, the capital ratio $x$ and the
treated share $\varpi$, by Propositions~\ref{prop:sp:B:circularity}
and~\ref{prop:sp:A:circularity}. Without policy, only the first margin responds: a rising material
price raises $P^T$ and hence $x$ and $a(x)$, but $\varpi$ does not move. Half of the transition
mechanism is switched off, and the half that survives is the one that operates through capital
deepening rather than through how much of the waste stream is captured.

Two other distortions are visible in~\eqref{eq:mkt:setup:lf-waste}. Waste is not treated for the sake
of the emissions treatment prevents --- the term $\left(1-d^W\right)\tau^E$ is missing --- and the mass
drawn into production is not priced at all, since $\tau^R=0$ makes $P^R=F_R$ rather than the full
social cost of a tonne. Both push the same way as the first, against circularity and towards
extraction. The gap between~\eqref{eq:mkt:setup:lf-waste} and its optimal counterpart is what the
instruments below close, and the second condition is exactly where the collection payment $s^W$ enters:
once the collector is paid per tonne collected, cost recovery no longer dictates how much it treats.


\subsection{Plan of Part~\ref{part:mkt}}\label{sec:mkt:setup:plan}

The two closures are developed separately and in the same order as in Part~\ref{part:sp}.
Section~\ref{sec:mkt:A} sets the instruments under Version~A, $\Phi^I=\Omega\phi^I$, and proves
implementation; Section~\ref{sec:mkt:B} does the same under Version~B, $\Phi^I$ exogenous, and
Section~\ref{sec:mkt:B:compare} collects the differences. Section~\ref{sec:mkt:schemes} then works
through the simplified waste-generation schemes of Section~\ref{sec:sp:schemes}, asking of each not
whether it changes the allocation --- Part~\ref{part:sp} settled that --- but whether it changes the
instrument set, and records the open items specific to the market economy.

Because the agents' problems have been stated once, the version sections are shorter than their
counterparts in Part~\ref{part:sp} and are not fully self-contained: each presupposes this section.
Each does, however, contain a complete statement of the equilibrium price system and of the instrument
values under its closure, so a reader who wants only one version needs Section~\ref{sec:mkt:setup} and
the section concerned.
